% OpenGauss / leanblueprint blueprint — "The non-Riemannian nature of perceptual
% color space, formalized." Feed to OpenGauss: /autoformalize blueprint/strangeness.tex
% \leanok = the Lean proof compiles (satreadout, Lean v4.28, CI green).
% Honesty convention: \begin{hypothesis} nodes are EMPIRICAL premises — they are
% NOT theorems and must never be \leanok. Everything downstream is conditional on them.

\section{The non-Riemannian nature of perceptual color space, formalized}

\paragraph{Thesis.} Perceived color difference is a \emph{saturating, subadditive}
readout, not a Riemannian (length) metric. Below we (i) fix the readout, (ii) prove
the algebra that forces diminishing returns, (iii) prove the geometric obstruction
— no midpoint, hence no length/Riemannian metric — and (iv) state honestly the one
\emph{empirical} link to human vision, which is not (and cannot be) formalized.

% ──────────────────────────────────────────────────────────────────────────────
\subsection{The empirical premise (not a theorem)}

\begin{hypothesis}[Saturating perception, Bujack et al.\ PNAS 2022]\label{hyp:bujack}
Human large-difference color judgments exhibit \emph{diminishing returns}: the
directly perceived difference between two colors is strictly less than the sum of
just-noticeable-difference steps between them. Operationally, perceived difference
is modelled as $D(c_1,c_2)=f_{\hat a}(\Delta E(c_1,c_2))$ with $f_a(t)=a(1-e^{-t/a})$
and a per-context saturation scale $\hat a$ fit from discrimination data
(e.g.\ $\hat a\approx 10$ on the $\Delta E_{00}$ anchor; $\hat a=0.1537$ in
OlmoEarth EO-embedding units). This is an empirical claim; it grounds, but is not
implied by, anything below.
\end{hypothesis}

% ──────────────────────────────────────────────────────────────────────────────
\subsection{The readout and its engine}

\begin{definition}[Saturating readout]\label{def:f}\lean{SatReadout.f}\leanok
For $a,t\in\mathbb{R}$, $f_a(t)=a\,(1-e^{-t/a})$. Local ruler ($f_a(0)=0$,
$f_a'(0)=1$, \lean{SatReadout.hasDerivAt_f_zero}) that bends to the asymptote $a$
(\lean{SatReadout.f_lt_asymptote}).
\end{definition}

\begin{lemma}[KEY identity]\label{lem:key}\lean{SatReadout.key_identity'}\uses{def:f}\leanok
$f_a(x)f_a(y)=a\bigl(f_a(x)+f_a(y)-f_a(x+y)\bigr)$ for all $a,x,y$ (only
\texttt{Real.exp\_add}). The subadditivity defect is exactly $f_a(x)f_a(y)/a$.
\end{lemma}

\begin{lemma}[Strict subadditivity = diminishing returns]\label{lem:subadd}
\lean{SatReadout.strict_subadd}\uses{lem:key}\leanok
$0<a,x,y \Rightarrow f_a(x+y)<f_a(x)+f_a(y)$. (Measured on the X11 gray ramp:
$\sum f(\text{step})=95.16$ vs $f(\text{direct})=10.00$, a $9.5\times$ overcount.)
\end{lemma}

\begin{remark}[The convexity route, and a mathlib gap]\label{rem:concave}
\uses{lem:subadd}
\Cref{lem:subadd} also follows from a general fact mathlib lacked:
\emph{a (strictly) concave $g$ on $[0,\infty)$ with $g(0)=0$ is (strictly)
subadditive}. We prove both forms,
\lean{ConcaveOn.subadditive_of_map_zero} and
\lean{StrictConcaveOn.subadditive_of_map_zero} (\leanok, $\mathbb{R}$ on
$[0,\infty)$), as a mathlib PR candidate (closest existing lemma:
\texttt{Convex\_subadditive\_le}, the reverse direction).
\end{remark}

% ──────────────────────────────────────────────────────────────────────────────
\subsection{The geometric obstruction (the strangeness)}

\begin{definition}[Saturated metric]\label{def:satdist}\lean{SatReadout.satDist}\uses{def:f}\leanok
$\mathrm{satDist}_a=f_a\circ\mathrm{dist}$ on a (pseudo)metric space; a pseudometric
(\lean{SatReadout.satPseudoMetricSpace}), with uniformity equal to the ambient one
(\lean{satPseudoMetricSpace_uniformity}).
\end{definition}

\begin{theorem}[No midpoint]\label{thm:nomid}\lean{SatReadout.no_midpoint}
\uses{lem:subadd,def:satdist}\leanok
In any metric space, for $p\ne q$ there is no $m$ with
$\mathrm{satDist}_a(p,m)=\mathrm{satDist}_a(m,q)=\tfrac12\mathrm{satDist}_a(p,q)$.
Quantitative $\varepsilon$-version: \lean{SatReadout.no_eps_midpoint} (gate
$2a\varepsilon<(D/2-\varepsilon)^2$).
\end{theorem}

\begin{theorem}[Collinear strictness in betweenness terms]\label{thm:wbtw}
\lean{SatReadout.collinear_strict_wbtw}\uses{lem:subadd}\leanok
Over a strict-convex normed $\mathbb{R}$-space: $\mathrm{Wbtw}\;\mathbb{R}\;p\,r\,q$
(via mathlib \texttt{dist\_add\_dist\_eq\_iff}) and $p\ne r\ne q$ imply
$\mathrm{satDist}_a(p,q)<\mathrm{satDist}_a(p,r)+\mathrm{satDist}_a(r,q)$: the
geodesic-between point is charged \emph{more} than the geodesic sum.
\end{theorem}

\begin{theorem}[Not a length / Riemannian metric]\label{thm:break}
\lean{SatReadout.satDist_breaks_geodesic}\uses{thm:nomid,thm:wbtw}\leanok
Over a strict-convex normed $\mathbb{R}$-space the base metric has the segment
midpoint (mathlib \texttt{eq\_midpoint\_of\_dist\_eq\_half}) while
$\mathrm{satDist}_a$ has none (\cref{thm:nomid}). A length space realises every
distance as an infimum of summed path segments and hence \emph{must} have
midpoints; therefore $f_a\circ\mathrm{dist}$ is not a length metric, a fortiori
not Riemannian.
\end{theorem}

\begin{theorem}[Not a geodesic metric — the formal non-Riemannian]\label{thm:notgeo}
\lean{satDist_not_geodesic}\uses{thm:nomid}\leanok
Call a metric \emph{geodesic} if every pair is joined by a metric segment
(\lean{IsGeodesicMetric}; ordinary $\mathbb{R}$ is, via $t\mapsto(1-t)p+tq$). For
any metric space with $\ge 2$ points and $0<a$, $\mathrm{satDist}_a$ is \emph{not}
geodesic: every geodesic metric is midpoint-convex
(\lean{IsGeodesicMetric.midpointConvex}), which \cref{thm:nomid} refutes. By
Hopf--Rinow a complete Riemannian metric is geodesic; hence $\mathrm{satDist}_a$
is the distance of no complete Riemannian structure.
\end{theorem}

\begin{theorem}[Not even a length metric — Hopf--Rinow-free]\label{thm:notlen}
\lean{satDist_not_length_metric}\uses{thm:nomid}\leanok
Call a metric \emph{intrinsic / length} if every pair has approximate midpoints for
every $\varepsilon>0$ (\lean{IsLengthMetric}). Then $\mathrm{satDist}_a$ is not a
length metric --- it has no approximate midpoints, by the quantitative
\lean{SatReadout.no_eps_midpoint}. This is strictly stronger than \cref{thm:notgeo}.
A Riemannian distance \emph{is} a length metric by definition, so this yields
non-Riemannian with \emph{no} Hopf--Rinow.
\end{theorem}

\begin{corollary}[Strangeness is monotone and bounded]\label{cor:strange}
\uses{thm:break}
Define strangeness of a pair as $\Delta E / f_a(\Delta E)$ (Riemannian vs.\ readout
verdict). It is $\approx 1$ for $\Delta E\lesssim$ a few JND (both metrics agree —
the regime where CIEDE is valid) and grows without bound as $\Delta E$ increases,
because $f_a\to a$. Worked over the X11 catalogue (676 colors): least strange =
\texttt{magenta} and its neighbourhood ($1.0\times$); most strange = the greens
(\texttt{green1}, $15.7\times$), i.e.\ magenta's opponent complement
(``$-$green''). The non-spectral hue and its complement bracket exactly the span
where the Riemannian ruler is most confidently wrong.
\end{corollary}

% ──────────────────────────────────────────────────────────────────────────────
\subsection{What is formal, and what is not}
The Lean development proves the \emph{geometry of the saturating readout}:
\cref{def:f}--\cref{cor:strange} are unconditional theorems about $f_a$ and
$\mathrm{satDist}_a$. Two boundaries, both named:
\begin{itemize}
\item \textbf{Empirical (not a theorem, by design).} The bridge to human vision is
\cref{hyp:bujack} alone — that perception \emph{is} a saturating readout. Data may
refine or refute it; it is never \leanok.
\item \textbf{Classical, now reduced to a definition.} \cref{thm:notlen}
machine-checks that $\mathrm{satDist}_a$ is \emph{not a length metric}. The only
step from there to ``Riemannian'' is the \emph{definitional} identity ``a Riemannian
distance is the infimum of path lengths, hence intrinsic/length'' --- not a theorem
to invoke but a definition. (The weaker \cref{thm:notgeo} would route through
Hopf--Rinow; \cref{thm:notlen} avoids it.)
\end{itemize}
Precise formalized statement: \emph{if} perceptual difference is a saturating
readout (\cref{hyp:bujack}), \emph{then} perceptual color space is not a length
metric (\cref{thm:notlen}, machine-checked), hence --- since Riemannian distance is
length by definition --- non-Riemannian.
